% !TEX program = xelatex
\documentclass[14pt, a4paper, oneside]{memoir}

\input{./packages.tex}
\input{./styles.tex}
\input{./commands.tex}

\title{Обоснование стратегии оптимизации Q-функции}
\begin{document}

\begin{titlingpage}
	\maketitle
\end{titlingpage}

\section{Правдоподобие смеси}

Допустим мы имеем независимые наблюдения $x_1, \dots, x_n$ полученные из смеси распределений.

Плотность смеси распределений имеет следующий вид:

\[f(x ~|~ \Psi) = \sum_{i=1}^k \omega_i \cdot f_i (x~|~\theta_i)\]

где $\Psi$ --- вектор параметров смеси $\omega_i, \theta_i$, $f_i$ --- $i$-ая компонента смеси.

Правдоподобием смеси распределений при наблюдаемых данных будет называться функция:

\[\mathcal{L}(x ~|~ \Psi) = \prod_{i=1}^n f(x_i ~|~ \Psi) = \prod_{i=1}^n \left(\sum_{j=1}^k \omega_j \cdot f_j(x_i ~|~ \theta_j)\right)\]

Для оценки параметров достаточно максимизировать правдоподобие. В силу монотонности логарифма можно оптимизировать логарифм правдоподобия:

\[\ln \mathcal{L}(x ~|~ \Psi) = \sum_{i=1}^n \ln \left(\sum_{j=1}^k \omega_j \cdot f_j(x_i~|~ \theta_j)\right)\]


Однако данную вещь неудобно оптимизировать из-за присутствия логарифма суммы, поэтому мы делаем дополнительное предположение, начиная работать с так называемыми "полными данными".

\section{Правдоподобие полных данных}

Пусть мы имеем независимые наблюдения $(x_i, z_i)$, где $x_i$ --- элементы выборки, а $z_i$ -- компонента к которой этот элемент принадлежит.
Функцией правдоподобия для полных данных, мы будем называть:

\[\mathcal{L}_c(x, z ~|~ \Psi) = \prod_{i=1}^n \mathbb{P} (x_i, z_i = j~|~ \Psi)\]

Совместную вероятность внутри произведения можно расписать по цепному правилу:

\[\mathbb{P} (x_i, z_i = j ~|~ \Psi) = \mathbb{P}(x_i ~|~ z_i = j, \Psi) \cdot \mathbb{P}(z_i = j ~|~ \Psi)\]

Понятно, что если $i$-ое наблюдение принадлежит $j$-ой компоненте, то $\mathbb{P}(z_i = j ~|~ \Psi) = \omega_j$, а вероятность $\mathbb{P}(x_i~|~ z_i=j, \Psi)$ описывается функцией плотности $f_j(x_i~|~ \theta_j)$.

Можно представить $z_i$ как вектор из нулей и одной единицы, стоящей на $j$-ом месте: $z_i = (z_{i1}, \dots, z_{ik})$. Тогда вероятности можно записать как:

\[\mathbb{P}(z_i = j ~|~ \Psi) = \prod_{i=1}^k \omega_j^{z_{ij}}\]
\[\mathbb{P}(x_i ~|~ z_i = j, \Psi) = \prod_{i=1}^k (f_j(x_i ~|~ \theta_j))^{z_{ij}}\]

Итоговая совместная плотность:

\[\mathbb{P}(x_i, z_i = j ~|~ \Psi) = \left(\prod_{i=1}^k \omega_j^{z_{ij}}\right) \left(\prod_{i=1}^k f_j(x_i~|~\theta_j)\right) = \prod_{i=1}^k (\omega_j \cdot f_j(x_i~|~ \theta_j))^{z_{ij}}\]

Поэтому правдоподобие полных данных равно:

\[\mathcal{L}_c(x, z ~|~ \Psi) = \prod_{i=1}^n \prod_{j=1}^k (\omega_j \cdot f_j(x_i~|~ \theta_j))^{z_{ij}}\]

Стандартно, работаем с логарифмом правдоподобия:

\[\ln \mathcal{L}_c (x, z ~|~ \Psi) = \sum_{i=1}^n \sum_{j=1}^k z_{ij} \cdot (\ln \omega_j + \ln f_j(x_i ~|~ \theta_j))\]


\section{Q-функция}

Однако в обычных условиях мы не знаем $z_{ij}$. Пусть $Z_{ij}$ --- случайные величины соответствующие $z_{ij}$. Тогда введём Q-функцию:

\[Q(\Psi ~|~ \Psi^{(k)}) = \mathbb{E}_{Z |  x, \Psi^{(k)}}[\ln \mathcal{L}_c(x, Z ~|~ \Psi)]\]

Оператор матожидания линеен, так что можем расписать:

\[Q(\Psi ~|~ \Psi^{(k)}) = \sum_{i=1}^n \sum_{j=1}^k \mathbb{E}[Z_{ij} ~|~ x, \Psi^{(k)}] \cdot (\ln \omega_j + \ln f_j(x_i ~|~ \theta_j))\]

$Z_{ij}$ --- бинарная случайная величина, поэтому ее матожидание равно:

\[\mathbb{E}[Z_{ij} ~|~ x, \Psi^{(k)}] = P(Z_{ij} = 1 ~|~ x_i, \Psi^{(k)})\]

По формуле Байеса эта величина равна:

\[P(Z_{ij} = 1 ~|~ x_i, \Psi^{(k)}) = \frac{\omega_j^{(k)} \cdot f_j(x_i~|~ \theta_j^{(k)})}{\sum_{j=1}^k \omega_j^{(k)} \cdot f_j(x_i~|~ \theta_j^{(k)})}\]

Обозначим эту величину за $\gamma_{ij}^{(k)}$, тогда:

\[Q(\Psi ~|~ \Psi^{(k)}) = \sum_{i=1}^n \sum_{j=1}^k \gamma_{ij}^{(k)} \cdot (\ln \omega_j + \ln f_j(x_i ~|~ \theta_j))\]

\section{Покомпонентная оптимизация Q-функции}

Имеем:

\[Q(\Psi ~|~ \Psi^{(k)}) = \sum_{i=1}^n \sum_{j=1}^k \gamma_{ij}^{(k)} \cdot (\ln \omega_j + \ln f_j(x_i ~|~ \theta_j))\]

Разделим на две суммы:

\[Q(\Psi ~|~ \Psi^{(k)}) =  \sum_{i=1}^n \sum_{j=1}^k \gamma_{ij}^{(k)} \cdot \ln \omega_j +  \sum_{i=1}^n \sum_{j=1}^k  \gamma_{ij}^{(k)} \ln f_j(x_i ~|~ \theta_j)\]

Понятно что первое и второе слагаемое можно оптимизировать независимо. Рассмотрим второе слагаемое и поменяем в нем порядок суммирования:

\[\sum_{i=1}^n \sum_{j=1}^k  \gamma_{ij}^{(k)} \ln f_j(x_i ~|~ \theta_j) = \sum_{j=1}^k \sum_{i=1}^n \gamma_{ij}^{(k)}\ln f_j(x_i ~|~ \theta_j)\]

Очевидно что если рассмотреть для каждого индекса $j$ то получаем функции независящие друг от друга.

\end{document}
